\documentclass[10pt]{amsart}
\usepackage[margin=0.95in]{geometry}
\usepackage{amsmath,amssymb,amsthm,hyperref}
\newtheorem{proposition}{Proposition}[section]
\newtheorem{lemma}[proposition]{Lemma}
\newtheorem{theorem}[proposition]{Theorem}
\newtheorem{corollary}[proposition]{Corollary}
\newcommand{\Ocal}{\mathcal O}
\newcommand{\Gcal}{\mathcal G}
\newcommand{\gfr}{\mathfrak g}
\newcommand{\C}{\mathbf C}
\newcommand{\Irr}{\operatorname{Irr}}
\newcommand{\Hom}{\operatorname{Hom}}
\newcommand{\val}{\operatorname{val}}
\title{Congruence data and compact induction in small characteristic}
\author{}
\date{Research note, 2 November 2026}
\begin{document}
\maketitle

This note gives congruence data that produce supercuspidal representations by compact
induction for every connected reductive group over every nonarchimedean local field.
Theorem~\ref{agc-all-groups} imposes no restriction on the residue characteristic or
the ramification of a splitting field. On a central extension with simply connected
derived group, elliptic toral functionals give deep congruence characters whose
intertwining is contained in an open subgroup compact modulo center. A finite
induction supplies the inducing representations. A smooth central twist then lets
the entire inducing pair descend to the original group.

The proof identifies group intertwining with transport of the full coadjoint coset,
using exact equivariant coordinates on unipotent radicals and a fixed depth margin
for the remaining analytic coordinates. Properness of the action on rationally
elliptic covectors supplies the required compact subgroup. The earlier sections
also compute root subgroups with different depths that restore a perfect finite
commutator pairing. For $\mathrm{Sp}_4$ over every local field of characteristic two,
we compute the inducing subgroup and all representations containing an explicit
family of characters, including characters whose leading reductions are not elliptic.

The Heisenberg and Clifford steps follow the architecture of Yu's construction and
its refinements \cite{Yu,FintzenSchwein,Takaya}. The occurrence statements for this
construction exhaust the representations containing the specified data. They do not supply
extraction from an arbitrary supercuspidal representation. The depth margin can
always be arranged for the constructed data; it cannot be imposed on a character
already known to occur without a further proof.

Existence of supercuspidal representations for all these groups is already known
by Kret's finite-group construction \cite{Kret}. The subject here is the specified
congruence construction and its intertwining calculation. Its outputs do not
exhaust all supercuspidals: Proposition~\ref{dzo-no-compact-transporter} and
Corollary~\ref{dzo-z-extensions} give an explicit depth-zero counterexample to
universal occurrence in this family, even after central extension and twisting.

Theorem~\ref{qtc-finite-equivalence} compares any two supplied compact inducing
pairs by finitely many matrix equations with an explicit Cartan bound. Thus
equivalence of successful inducing data is already controlled. For extraction,
Proposition~\ref{stre-finite-refinement} gives a finite refinement that preserves
occurrence under an explicit precision hypothesis. The remaining coordinates
retain the radical of the commutator pairing. Their compatibility with earlier
stages is not established. Propositions~\ref{mkdo-annular-types} and
\ref{mkdo-fixed-vertex} also show why increasing the level at a fixed vertex
cannot serve as a universal extraction procedure.

Proposition~\ref{ofd-multiplicity-saturation} gives the exact formal-degree
criterion for an occurring representation of an open subgroup compact modulo
center to induce the prescribed supercuspidal. It also distinguishes an
irreducible compact induction from a compact induction consisting of several
copies of one supercuspidal. The latter has a matrix algebra of endomorphisms;
choosing a primitive idempotent in that algebra does not by itself produce a
new inducing subgroup.

Theorem~\ref{frc-bounded-decision} makes closure under the coefficient blocks
a finite decision for every supplied irreducible local model. A radius
depending on the local subgroup and fiber dimension bounds every possible
finite block. The procedure either finds inducing data or certifies failure
of that model. Theorem~\ref{svpe-extraction} gives a positive extraction
criterion: an occurring type at a full special vertex induces the given
supercuspidal if its restriction to the intersection with a minimal
parabolic is irreducible. Neither result asserts that every supercuspidal
has a local model satisfying its successful alternative.

Theorem~\ref{fel-polynomial-locus} also makes a change of embedding within
a fixed finite multiplicity space a finite system of polynomial equations.
Proposition~\ref{fel-empty-locus} shows that this system can have no solution
even when the representation is already compactly induced. Its proof gives
an explicit change of local type that succeeds in that example. Thus
changing the embedding alone does not provide universal extraction.

Theorem~\ref{girg-rank-gap} computes the global complement of the outgoing
coefficient blocks. Their total rank satisfies $R=n$ or $R\geq2n$ when
the coefficient projection is nontrivial. Thus the previously considered
strict-descent interval $n<R<2n$ is empty. The case $R=n$ gives compact
induction directly. Proposition~\ref{ecbe-endpoints} also extracts an
inducing pair at either endpoint of the diagonal bounds for every $R$.
At $R=2n$, Theorem~\ref{rcx-rotation} gives an explicit successful rotation
when $c=1/2$ and the complement has the original local type. These results
apply in every residue characteristic and impose no splitting hypothesis.
Theorems~\ref{toue-extraction} and \ref{toue-self-reversing} remove that
type hypothesis whenever the active coefficient blocks have two neighbor
orbits at $R=2n$. They construct equivariant unitary operators directly
from the blocks. Lemma~\ref{scs-saturation} identifies reducing subspaces
forced by equality in the rank inequality for the projection equations.
It yields actual compact induction at $R=2n$ if every range projection
has full central support in the local range algebra, if every individual
block has rank at most two, or if the fiber dimension is prime
(Theorems~\ref{scs-full-central}, \ref{scr-rank-two}, and
Corollary~\ref{scr-prime-fiber}). In these cases the closure of an original
fiber is supported on at most $4n$ coordinates. Theorem~\ref{scs-degree-three}
also handles a single active orbit of degree three.
Theorem~\ref{cra-extraction} treats every commutative local range algebra,
with arbitrary multiplicities, by an explicit decomposition into triangles
or four-cycles. Theorems~\ref{epc-spanning} and \ref{epc-probe-lattice} give
further criteria for noncommutative algebras acting with multiplicity one.
They recover central incidence spaces from the images of other range
projections, allowing finite sums and intersections. This permits arbitrarily
large individual block ranks. Corollary~\ref{fmd-prime-or-four} also handles
the spanning condition when the local factors have prime multiplicity or
multiplicity four. Corollary~\ref{lfr-extraction} completes extraction from
every boundary model with fiber dimension at most eleven.
These assertions concern supplied
coefficient models; they do not produce such models for every supercuspidal
representation.

Proposition~\ref{god-orbital-moves} gives further changes of local data from
individual coefficient orbits, without a bound on the total rank $R$.
Theorem~\ref{god-termination} bounds the number of steps that strictly
decrease the formal-degree ratio and then the fiber dimension. The bound
uses only the initial fiber dimension and compact subgroup volumes.
Corollary~\ref{god-terminal-core} extracts every terminal model with
$R\leq2n$. Models with larger rank can remain, so this finite descent
does not yet give universal exhaustion.

Theorem~\ref{wpc-exhaustion} and Corollary~\ref{wpc-actual-data} give an
exhaustive construction under the single condition $p\nmid|W|$, where $W$
is the absolute Weyl group. They require no tame splitting hypothesis on
$G$ over $F$. This consequence of the known exhaustion and equivalence
theorems uses Yu data over the fields defining the absolute simple factors,
a compatible central character, and finite Clifford theory. It includes
wildly ramified restriction of scalars and arbitrary central tori.

\section{The integral obstruction to changing a representative}

Let $\Ocal$ be a discrete valuation ring with uniformizer $\varpi$, fraction field $F$,
and normalized valuation $\val$. Let $L$ be a finite free $\Ocal$-module, and let
$R\subseteq L$ be a submodule of rank $s$. Write
$R^{\mathrm{sat}}=(FR)\cap L$. Choose Smith bases so that $L$ has basis
$e_1,\ldots,e_l$ and $R$ has basis
$r_i=\varpi^{a_i}e_i$, where $a_i\geq0$ and $1\leq i\leq s$.

\begin{lemma}\label{lem:smith-correction}
For $\lambda\in L^\vee$ and an integer $q\geq1$, the coset
$\lambda+\varpi^qL^\vee$ contains a functional vanishing on $R$ if and only if
\[
 \lambda(r_i)\in\varpi^{q+a_i}\Ocal
 \quad(1\leq i\leq s).
\]
Equivalently, $\lambda(R^{\mathrm{sat}})\subseteq\varpi^q\Ocal$.
If $\delta=\max_i a_i$, the stronger condition
$\lambda(R)\subseteq\varpi^{q+\delta}\Ocal$ suffices.
\end{lemma}

\begin{proof}
Vanishing on $r_i$ is equivalent to vanishing on $e_i$, since $\Ocal$ is a domain.
A correction in $\varpi^qL^\vee$ can change $\lambda(e_i)$ to zero exactly when
$\lambda(e_i)\in\varpi^q\Ocal$. The coordinates of the correction on the remaining
basis vectors can be zero. Since
$\lambda(r_i)=\varpi^{a_i}\lambda(e_i)$, this proves every assertion.
\end{proof}

Thus divisibility on the original coroot lattice is insufficient when that lattice
is not saturated. A bounded denominator gives a sufficient margin, but it does not
make that margin hold in a prescribed character coset. In equal characteristic $p$,
the differential of a coroot divisible by $p$ can vanish. Dividing the abstract
cocharacter by $p$ then changes the actual differential and the group commutator;
it is not an integral correction of a representative.

\section{A finite refinement using actual root subgroups}

Let $F$ be a nonarchimedean local field with ring of integers $\Ocal$, residue field
$k$ of cardinality $q$, and uniformizer $\varpi$. Fix an additive character
$\psi:F\to\C^\times$ whose annihilator of $\Ocal$ is $\Ocal$.
Let $\Gcal$ be a split connected reductive group scheme over $\Ocal$, with split
maximal torus $\mathcal T$ and Chevalley root homomorphisms $x_\alpha$.
Put $G=\Gcal(F)$ and $L=\operatorname{Lie}(\Gcal)(\Ocal)$. Use the decomposition
\[
 L=\operatorname{Lie}(\mathcal T)(\Ocal)
       \oplus\bigoplus_{\alpha\in\Phi}\Ocal E_\alpha,
 \qquad E_\alpha=dx_\alpha(1).
\]
Choose signs so that $[E_\alpha,E_{-\alpha}]=H_\alpha$, where
$H_\alpha=d\alpha^\vee(1)$. The signs in the alternating pairings below can be
changed by changing root coordinates.

For an integer $j\geq1$, let
$G_j=\ker(\Gcal(\Ocal)\to\Gcal(\Ocal/\varpi^j\Ocal))$.
Choose $\beta\in L^\vee$ vanishing on all root spaces. Define
\[
 c_\alpha=\beta(H_\alpha),\qquad
 \Gamma=\{\alpha\in\Phi:c_\alpha\ne0\},\qquad
 d_\alpha=\val(c_\alpha)\quad(\alpha\in\Gamma).
\]
Then $d_{-\alpha}=d_\alpha$. Set $\delta=\max_{\alpha\in\Gamma}d_\alpha$,
with $\delta=0$ if $\Gamma$ is empty. Choose an integer $m\geq2\delta+1$, and put
\[
 n=2m+1,\qquad H=G_{m+1},\qquad a_\alpha=m-d_\alpha.
\]
Since $2(m+1)\geq n$, the square-zero congruence quotient gives a canonical map
\[
 \ell:H/G_n\xrightarrow{\sim}
       \varpi^{m+1}L/\varpi^nL.
\]
Define the character $\theta$ of $H$ by
$\theta(h)=\psi(\varpi^{-n}\beta(\ell(h)))$. This is independent of the lift of
$\ell(h)$ to $\varpi^{m+1}L$.
Finally, let $P$ be the subgroup generated by $H$ and
$x_\alpha(\varpi^{a_\alpha}\Ocal)$ for $\alpha\in\Gamma$.

\begin{proposition}\label{prop:weighted-heisenberg}
The group $P$ is compact and open, $H$ is normal in $P$, and $\theta$ is
$P$-invariant. Root coordinates identify its abelian quotient with
\[
 V=P/H\simeq\bigoplus_{\alpha\in\Gamma}\Ocal/\varpi^{d_\alpha+1}\Ocal.
\]
The commutator pairing $B_\theta$ on $V$ is trivial between nonopposite root
summands. On a pair of opposite summands it is
\begin{equation}\label{eq:root-pairing}
 B_\theta\bigl(x_\alpha(\varpi^{a_\alpha}u),
                    x_{-\alpha}(\varpi^{a_\alpha}v)\bigr)
 =\psi\left(\varpi^{-(d_\alpha+1)}
                   (\varpi^{-d_\alpha}c_\alpha)uv\right).
\end{equation}
In particular, this pairing is perfect. There is a unique irreducible smooth
representation $\omega_\theta$ of $P$ whose restriction to $H$ is
$\theta$-isotypic, and
\[
 \dim\omega_\theta
   =\prod_{\alpha\in\Gamma\cap\Phi^+}q^{d_\alpha+1}.
\]
Let $C$ be the inverse image in $P$ of the sum of the positive-root summands
of $V$. The character $\theta$ extends to a character $\eta$ of $C$, and
$\omega_\theta\simeq\operatorname{Ind}_C^P\eta$ for every such extension.
For every smooth representation $\pi$ of $G$, evaluation gives an isomorphism
of $P$-representations
\begin{equation}\label{eq:occurrence}
 \omega_\theta\otimes
       \Hom_P(\omega_\theta,\pi)
       \xrightarrow{\sim}\pi^\theta.
\end{equation}
Consequently $\theta$ occurs in $\pi$ if and only if $\omega_\theta$ occurs in $\pi$.
Moreover,
\begin{equation}\label{eq:intertwining-inclusion}
 I_G(\omega_\theta)\subseteq I_G(\theta).
\end{equation}
\end{proposition}

\begin{proof}
Put $a=m-\delta\geq1$. Since $2a\geq m+1$, the quotient
$G_a/G_{m+1}$ is an additive congruence quotient. The images of the indicated
root subgroups are its independent root submodules, with the indicated lengths.
This proves the description of $P/H$. It also proves compactness and openness:
$P$ is a union of finitely many cosets of $H$ in $\Gcal(\Ocal)$.
The group $H$ is normal in $\Gcal(\Ocal)$.

For a toral functional, the integral adjoint action of a root element satisfies
\[
 \beta\circ\operatorname{Ad}(x_\alpha(t))-\beta
       =\pm t c_\alpha E_{-\alpha}^\vee.
\]
Indeed, the only root space that can acquire a toral component under this
action is the $-\alpha$ root space, and all subsequent terms lie in root spaces.
Conjugation acts on the square-zero congruence quotient by this integral adjoint
action. For $t\in\varpi^{a_\alpha}\Ocal$, the change in its pairing with an
element of $\varpi^{m+1}L$ has valuation at least
$a_\alpha+d_\alpha+m+1=n$. Thus each root generator preserves $\theta$.
Invariance under $H$ follows because $\theta$ is a character.

The integral Chevalley commutator formula for nonopposite roots is a product
of root elements with parameter valuations at least
$a_\alpha+a_\gamma\geq2m-2\delta\geq m+1$. Each factor belongs to $H$, and
$\theta$ is trivial on every root subgroup in $H$. Their pairing is therefore
trivial.

For opposite roots use the root homomorphism from $\mathrm{SL}_2$.
Write $u'=\varpi^{a_\alpha}u$, $v'=\varpi^{a_\alpha}v$, and $s=u'v'$.
In $\mathrm{SL}_2$ the commutator is
\[
 \begin{pmatrix}1&u'\\0&1\end{pmatrix}
 \begin{pmatrix}1&0\\v'&1\end{pmatrix}
 \begin{pmatrix}1&-u'\\0&1\end{pmatrix}
 \begin{pmatrix}1&0\\-v'&1\end{pmatrix}
 =
 \begin{pmatrix}
  1+s+s^2&-(u')^2v'\\
  u'(v')^2&1-s
 \end{pmatrix}.
\]
Its lower-diagonal-upper factorization has torus
$\operatorname{diag}(A,A^{-1})$, with $A=1+s+s^2$, and root parameters
$u'(v')^2/A$ and $-(u')^2v'/A$. Those root parameters have valuation at least
$3a_\alpha\geq m+1$, so their images lie in $H$ and have trivial $\theta$-value.
The remaining factor is $\alpha^\vee(A)$. In its congruence linearization,
\[
 \ell(\alpha^\vee(A))=u'v'H_\alpha
       \pmod{\varpi^{4a_\alpha}L}.
\]
Here both the term $s^2$ and the nonlinear terms of the integral cocharacter
have valuation at least $4a_\alpha$. The bound on $m$ implies
$4a_\alpha\geq4m-4\delta\geq n$. Applying $\theta$ therefore gives
$\psi(\varpi^{-n}c_\alpha u'v')$, which is \eqref{eq:root-pairing}.
Because $\varpi^{-d_\alpha}c_\alpha$ is a unit and the annihilator of
$\Ocal$ under $\psi$ is $\Ocal$, each opposite-root pairing is perfect.

Let $N=\ker\theta$. Invariance makes $N$ normal in $P$, and $P/N$ is finite.
The subgroup $H/N$ is central, with faithful character induced by $\theta$.
Perfection of the commutator pairing makes $H/N$ the whole center of $P/N$.
The central-character summand of $\C[P/N]$ has dimension $|P/H|$ and center
$\C$: the coefficient of a noncentral coset must vanish in any central element,
since some commutator is detected by $\theta$. It is therefore a single matrix
algebra. This proves uniqueness and the dimension formula.

The positive-root summands form a maximal isotropic subgroup of $V$.
Consequently $C/N$ is abelian, so its central subgroup $H/N$ has an extension
of $\theta$ to $C/N$. Inflate an extension to $C$ and call it $\eta$.
The induced representation $\operatorname{Ind}_C^P\eta$ is
$\theta$-isotypic on $H$ and has dimension $[P:C]=\sqrt{|P/H|}$.
Uniqueness and the dimension formula just proved show that it is
$\omega_\theta$.

The space $\pi^\theta$ is $P$-stable and is a module for that matrix algebra.
Evaluation proves \eqref{eq:occurrence}, including for an infinite-dimensional
smooth $\pi$. Finally, a nonzero intertwiner for $\omega_\theta$ restricts to
$H\cap gHg^{-1}$, where the two representations act by the scalar characters
$\theta$ and ${}^g\theta$. These characters must agree, proving
\eqref{eq:intertwining-inclusion}.
\end{proof}

\section{A characteristic-two example in type $C_2$}

Take $F=k((t))$ of characteristic two, $\Gcal=\mathrm{Sp}_4$ over $k[[t]]$, and
the diagonal torus with coordinates
$\operatorname{diag}(z_1,z_2,z_1^{-1},z_2^{-1})$.
Let $h_i$ be the differential of its $i$th coordinate cocharacter. Choose
$\beta(h_1)=1$ and $\beta(h_2)=1+t$. The long roots $2e_i$ have coroots $h_i$,
whereas the short roots $e_1-e_2$ and $e_1+e_2$ have coroots $h_1-h_2$ and
$h_1+h_2$. Thus
\[
 d_\alpha=0\text{ for long roots},\qquad
 d_\alpha=1\text{ for short roots}.
\]
For every $m\geq3$, put $H=G_{m+1}$ and $n=2m+1$. The group $P$ is generated
by $H$, the short root subgroups at depth $m-1$, and the long root subgroups
at depth $m$. Proposition~\ref{prop:weighted-heisenberg} gives
\[
 P/H\simeq(\Ocal/t^2\Ocal)^4\oplus k^4,
 \qquad \dim\omega_\theta=q^6.
\]
All commutator assertions use the actual group. In particular, nonopposite
short-root commutators that have a long-root term with coefficient two vanish
in characteristic two. Mixed short-long commutators are already in $H$.
For opposite short roots, the possible nonlinear torus error has depth at least
$4(m-1)\geq2m+1$; for long roots it has depth at least $4m$.

For $k=\mathbf F_2$, the group $P/\ker\theta$ is an extraspecial group of order
$2^{13}$. Its center is $H/\ker\theta$ of order two, and its quotient by that
center is the twelve-dimensional $\mathbf F_2$-space $V$. Define a quadratic
form $Q:V\to\mathbf F_2$ by
$(-1)^{Q(v)}=\theta(\widetilde v^2)$ for a lift $\widetilde v\in P$.
Each root element has order two. The positive-root half of $V$ is isotropic
for the commutator pairing and is generated by these square-one lifts.
It is therefore a totally singular subspace of dimension six for $Q$.
Thus $Q$ has Arf invariant zero. The preimage of this half in
$P/\ker\theta$ is elementary abelian; choose the extending character $\eta$
with values in $\{1,-1\}$. Its induced representation gives a real model of
the irreducible Heisenberg representation of dimension $64$. This assertion
does not include an extension of that representation to a normalizer.

The roots on which the leading reduction of $\beta$ vanishes are precisely the
short roots. They form a root system of type $D_2$, which is not the root system
of a proper Levi subgroup of $\mathrm{Sp}_4$. Nevertheless the refinement above
has a perfect commutator pairing and preserves occurrence. This is a statement
about this finite step, not an exclusion of possible data with several stages.
No connectedness assertion about the full coadjoint stabilizer is used. In
characteristic two the Weyl sign changes act trivially on the toral Lie algebra
and can give additional stabilizer components.

The same valuation calculation should not be transferred unchanged to mixed
characteristic. Over $\mathbf Q_2$ with $\varpi=2$, the values $1$ and $1+\varpi=3$
give short-root values $-2$ and $4$, of valuations one and two. Here the displayed
general bound is $m\geq5$ and the Heisenberg dimension is $q^7$.

\input{split-hyperspecial-intertwining-section}
\input{elliptic-congruence-sector-section}
\input{all-groups-deep-intertwining-section}
\input{rational-elliptic-properness-section}
\input{toral-ellipticity-section}
\input{central-twist-descent-section}
\input{all-group-construction-section}
\input{depth-zero-obstruction-section}
\input{maximal-compact-depth-obstruction-section}
\input{horospherical-compression-obstruction-section}
\input{finite-reducing-closure-decision-section}
\input{finite-embedding-locus-section}
\input{rank-budget-extraction-section}
\input{global-incidence-rank-gap-section}
\input{rotation-complement-extraction-section}
\input{two-orbit-unitary-extraction-section}
\input{saturation-central-support-section}
\input{rank-two-central-refinement-section}
\input{commuting-range-algebra-extraction-section}
\input{exclusive-probe-central-refinement-section}
\input{factor-multiplicity-descent-section}
\input{low-fiber-rank-boundary-corollary}
\input{general-orbital-descent-section}
\input{small-characteristic-transverse-extraction-section}
\input{rational-cocharacter-levi-section}
\input{quantitative-transporter-comparison-section}
\input{occurrence-formal-degree-section}
\input{special-vertex-parabolic-extraction-section}
\input{weyl-prime-compact-induction-section}

\section{The remaining extraction problem}

For $G=\mathrm{PGL}_2$ in mixed characteristic with residue characteristic two,
write $h$ for the class of $\operatorname{diag}(1,0)$ in its integral toral Lie
algebra. The coroot differential is $2h$. If $\beta(h)=1$, the coroot-lattice
Smith exponent is $e=\val(2)$. Proposition~\ref{prop:weighted-heisenberg} applies
with $m\geq2e+1$, giving
\[
 P/H\simeq(\Ocal/\varpi^{e+1}\Ocal)^2,
 \qquad\dim\omega_\theta=q^{e+1}.
\]
This repairs the finite pairing by changing the actual root subgroups.

It does not give a generic representative in the original filtration coset.
Identify the dual of $\mathfrak{pgl}_2(\Ocal)$ with trace-zero matrices by the trace
pairing, and take $B=\operatorname{diag}(1,-1)$. Every matrix
$B'\in B+\varpi\mathfrak{sl}_2(\Ocal)$ has unit determinant and eigenvalues
$\lambda,-\lambda$ of valuation zero over a splitting field. Its coroot value is
$2\lambda$, of valuation $e$, whereas its depth at the original vertex is zero.
Thus no such representative satisfies the toral genericity equality at depth
zero. Scaling the coset shifts the two valuations equally and leaves their gap
unchanged. The correction margin in Lemma~\ref{lem:smith-correction} cannot be
discarded.

In equal characteristic two the same coroot differential is zero. Replacing it
by the formally divided vector $h$ does not change the zero differential of the
actual root commutator. More generally, for a central homomorphism
$f:\widetilde G\to G$ and a pulled-back character, one has the exact identity
\[
 B_{\theta\circ f}(u,v)=B_\theta(f(u),f(v))
\]
whenever the indicated finite commutator quotients are defined. A pulled-back
pairing cannot acquire directions absent from the actual commutator by a change
of terminology. Changing filtrations or retaining additional central and
Clifford data is a further operation that needs its own occurrence argument.

Propositions~\ref{prop:ecs-sp4-sector} and \ref{prop:ecs-inertia} supply compactness,
the actual inducing data, and exhaustion of the specified characteristic-two
sectors. Their proof uses the entire character coset and the precision bound
$m\ge2\delta+1$. An arbitrary occurring minimal-depth character need not satisfy
that bound. Theorem~\ref{agc-all-groups} removes the group and splitting restrictions
for a nonempty construction. Its central twist is applied after constructing a
representation and can change both depth and the frequency of a congruence character.
It therefore does not assert a correction within the original character coset.

Proposition~\ref{dzo-no-compact-transporter} proves that no choice of the single
deep linear character in Theorem~\ref{agc-all-groups} can cover the representation
\eqref{dzo-representation}; Corollary~\ref{dzo-z-extensions} excludes repair by its
central extensions and smooth twists. Thus universal occurrence in this family
is false. The depth-zero cuspidal datum already gives the correct compact induction
for that example. A full classification must allow a larger collection of data,
including cuspidal representations of finite reductive quotients and further stages
at positive depth. Extracting such data while preserving occurrence, and proving
that extraction terminates for every supercuspidal, remain unresolved here.

One sufficient relative step would retain a smaller twisted Levi group $M$,
a compact inducing representation $\tau$ on $K_M\subset M(F)$, and an
occurring representation $\lambda$ on a subgroup $K=JK_M\subset G(F)$,
open and compact modulo $Z(G)(F)$. Here $J$ denotes the additional compact
subgroup supplied by the step. The essential intertwining calculation would be
\[
 I_{G(F)}(\lambda)\subseteq J I_{M(F)}(\tau)J.
\]
If the lower compact induction is irreducible, so that
$I_{M(F)}(\tau)=K_M$, this inclusion confines all intertwining to $K$.
For irreducible $\lambda$, the compact-induction argument then identifies
the given representation with $\operatorname{c-Ind}_K^{G(F)}\lambda$.
Depth-zero exhaustion supplies a terminal case. What is missing is extraction
of the occurring data and this relative intertwining comparison in small
characteristic. The existence of a suitable reductive $M$ is itself part
of that problem; the coordinate calculations above do not establish it.

Proposition~\ref{stre-finite-refinement} supplies a genuine local part of
extraction: it chooses a character that occurs, then corrects it by an actual
group element. Lemma~\ref{stre-schur-complement} computes the exact residual
commutator form, including its quadratic correction. Neither result forces the
precision hypothesis for an arbitrary input or gives a compatible sequence
of residual groups and representations. Proposition~\ref{stre-pgl2-closed-orbits}
shows that geometric closedness of the orbit cannot be imposed on every
occurring frequency coset in characteristic two.

Proposition~\ref{rcl-levi-characterization} describes rationally
cocharacter-closed covectors using elliptic covectors on $F$-Levi subgroups.
The description holds in every characteristic. The Levi subgroup in this
description need not be the stabilizer of the covector. In
Proposition~\ref{rcl-non-pseudoreductive-stabilizer}, the stabilizer is
noncompact modulo center and contains a nontrivial smooth connected
unipotent normal subgroup. Thus replacing geometric closedness by rational
cocharacter closure requires a further representation-theoretic descent
argument; the orbit description does not supply one.

The finite equivalence theorem applies once inducing pairs have been obtained,
including pairs arising through different central extensions and twists.
For a unitary supercuspidal $\pi$, Proposition~\ref{ofd-multiplicity-saturation}
expresses the extraction target as an occurring pair $(J,\tau)$ satisfying
\[
 \dim\tau=d_\pi\operatorname{vol}(J/Z).
\]
Every occurring pair satisfies the opposite weak inequality. Proving that
equality occurs somewhere is an existence assertion; it does not follow from
the inequality or from finite-dimensionality of an endomorphism algebra.
Proposition~\ref{hco-positive-defect} also tests a geometric compression
on an actual supercuspidal representation. Removing the maximally descending
Busemann blocks leaves the nonzero defect
$QPQ-(QPQ)^2=\tfrac29Q$. A geometric extraction must produce a subspace
that reduces the coefficient projection, rather than only annihilating its
extreme matrix blocks.

Theorem~\ref{frc-bounded-decision} bounds the closure search for an arbitrary
open subgroup compact modulo center. It terminates even when
$Z(\mathcal A)=0$: in that case it returns a finite certificate of failure
for the supplied coordinate model. The representation can still have
compact inducing data in a different model. Theorem~\ref{svpe-extraction}
proves a positive result under a finite restriction condition on an
occurring special-vertex type. Neither the existence of such a type for an
arbitrary input nor a change of model guaranteed to succeed is established.

For a fixed occurring local type, Theorem~\ref{fel-polynomial-locus}
replaces the search over embeddings by finitely many equations.
Proposition~\ref{fel-two-points} gives a full multiplicity space whose
successful locus consists of two points. Proposition~\ref{fel-empty-locus}
gives another occurring type whose entire locus is empty, although its
supercuspidal representation has an explicit edge inducing pair. The
compatibility criterion in Proposition~\ref{fel-compatible-induction}
explains the distinction: the chosen local type must be a full induction
from its intersection with an inducing subgroup, not just a constituent
of that induction. A universal procedure must therefore change local
types as well as embeddings and prove that a successful type is reached.

The global complement identity in Theorem~\ref{girg-rank-gap} shows that
$n<R<2n$ is impossible, and every model with $R<2n$ already extracts.
Both endpoints of the diagonal bounds give finite coordinate blocks by
Proposition~\ref{ecbe-endpoints}. At $R=2n$, a reducible complement gives
a smaller occurring type. Theorem~\ref{rcx-rotation} instead changes the
embedding explicitly when $c=1/2$ and the complement is isomorphic to the
original type. The rank of the coefficient blocks after a general change
of type is not controlled by these statements. At the same rank boundary,
Theorems~\ref{toue-extraction} and \ref{toue-self-reversing} handle every
model with two active neighbor orbits. The equality argument in
Lemma~\ref{scs-saturation} gives further extractions within the original
coefficient model. Theorems~\ref{scs-full-central}, \ref{scs-degree-three},
and \ref{scr-rank-two}, together with Corollary~\ref{scr-prime-fiber},
handle full central support, degree three, individual block ranks at most
two, and prime fiber dimensions. Theorem~\ref{cra-extraction} handles every
commutative local range algebra. Theorem~\ref{epc-spanning} also treats
noncommutative algebras acting with multiplicity one under the incidence
spanning condition \eqref{epc-one-factor-spanning}. Consequently an
uncovered boundary model has a single
orbit preserved by reversal, degree at least four, individual block rank
at least three, and composite fiber dimension at least twelve, by
Corollary~\ref{lfr-extraction}. Its local range algebra is noncommutative and
has more than one central summand, some of which are omitted by an active
range projection. In multiplicity one it also fails the more general
recovery condition of Theorem~\ref{epc-probe-lattice}. Under the spanning
condition, Corollary~\ref{fmd-prime-or-four} further excludes prime
multiplicities and multiplicity four. These restrictions do not exclude
every remaining local algebra.

The individual-orbit moves of Proposition~\ref{god-orbital-moves} are
defined for every total rank. Their strict lexicographic descent terminates
within the explicit bound of Theorem~\ref{god-termination}. A terminal
model with $R\leq2n$ gives compact induction, but a terminal model with
$R>2n$ is not excluded. Corollary~\ref{god-minimum-pair} proves the same
necessary rank inequality for a pair minimizing the ratio and then the
fiber dimension. The finite descent is not shown to attain that global
minimum or to produce an inducing pair for every representation.

When $p\nmid|W|$, Theorem~\ref{wpc-exhaustion} proves exhaustion without
assuming tame splitting over $F$, and Corollary~\ref{wpc-actual-data}
specifies the inducing data. Its proof uses the full-pair conjugacy theorem
for Yu data to control the finite central-isogeny quotient. At primes
dividing $|W|$, neither exhaustion on the absolute factors nor this
conjugacy mechanism is supplied by that argument.

The unresolved assertion is universal extraction of such a pair. Even the
underlying compact-induction statement remains Conjecture~1.1 in
\cite{JohnsonWeissman}; that work proves a geometric extraction theorem in
relative rank one and formulates higher-rank conjectures. The local results
here do not establish unrestricted higher-rank extraction or classification.

\begin{thebibliography}{99}
\bibitem{Yu}
J.-K. Yu, \emph{Construction of tame supercuspidal representations},
J. Amer. Math. Soc. \textbf{14} (2001), 579--622.
\bibitem{FintzenSchwein}
J. Fintzen and D. Schwein, \emph{Construction of tame supercuspidal representations
in arbitrary residue characteristic}, arXiv:2501.18553,
\url{https://arxiv.org/abs/2501.18553}.
\bibitem{Takaya}
Y. Takaya, \emph{Characteristic-free approaches around Yu's construction},
arXiv:2605.03638v2, \url{https://arxiv.org/abs/2605.03638}.
\bibitem{HR}
T. J. Haines and S. Rostami, \emph{The Satake isomorphism for special maximal
parahoric Hecke algebras}, Represent. Theory \textbf{14} (2010), 264--284,
\url{https://arxiv.org/abs/0907.4506}.
\bibitem{CF}
B. Conrad, lectures, and T. Feng, notes, \emph{Reductive groups over fields},
version of 5 December 2020,
\url{https://math.stanford.edu/~conrad/249BW16Page/handouts/AMSalggroups2.pdf}.
\bibitem{Cotner}
S. Cotner, \emph{Exercises for Arizona Winter School 2025},
\url{https://public.websites.umich.edu/~stcotner/Expository/AWS_2025.pdf}.
\bibitem{KP}
T. Kaletha and G. Prasad, \emph{Bruhat--Tits theory---a new approach},
New Mathematical Monographs \textbf{44}, Cambridge University Press, 2023.
\bibitem{Kret}
A. Kret, \emph{Existence of cuspidal representations of $p$-adic reductive groups},
arXiv:1205.2771, \url{https://arxiv.org/abs/1205.2771}.
\bibitem{MP96}
A. Moy and G. Prasad, \emph{Jacquet functors and unrefined minimal $K$-types},
Comment. Math. Helv. \textbf{71} (1996), 98--121,
\url{https://doi.org/10.1007/BF02566411}.
\bibitem{FintzenTypes}
J. Fintzen, \emph{Types for tame $p$-adic groups},
Ann. of Math. (2) \textbf{193} (2021), 303--346,
\url{https://doi.org/10.4007/annals.2021.193.1.4}.
\bibitem{FintzenTori}
J. Fintzen, \emph{Tame tori in $p$-adic groups and good semisimple elements},
Int. Math. Res. Not. IMRN 2021, no. 19, 14882--14904,
\url{https://doi.org/10.1093/imrn/rnz234}.
\bibitem{HakimMurnaghan}
J. Hakim and F. Murnaghan, \emph{Distinguished tame supercuspidal representations},
Int. Math. Res. Pap. IMRP 2008, rpn005,
\url{https://doi.org/10.1093/imrp/rpn005}.
\bibitem{KalethaRegular}
T. Kaletha, \emph{Regular supercuspidal representations},
J. Amer. Math. Soc. \textbf{32} (2019), 1071--1170,
\url{https://doi.org/10.1090/jams/925}.
\bibitem{JohnsonWeissman}
S. Johnson and M. H. Weissman, \emph{Types and collapse for cuspidal
representations of groups acting on trees}, arXiv:2607.27590,
\url{https://arxiv.org/abs/2607.27590}.
\bibitem{BHMR}
M. Bate, S. Herpel, B. Martin, and G. R\"ohrle,
\emph{Cocharacter-closure and the rational Hilbert--Mumford Theorem},
Math. Z. \textbf{287} (2017), 39--72,
\url{https://doi.org/10.1007/s00209-016-1816-5}.
\end{thebibliography}

\end{document}
